% Generated draft: requires mathematical and source review.
\documentclass[11pt]{article}
\usepackage[margin=25mm]{geometry}
\usepackage{fontspec}
\setmainfont{Latin Modern Roman}
\usepackage{amsmath,amssymb,mathtools,unicode-math}
\setmathfont{Latin Modern Math}
\usepackage{enumitem,needspace,xurl,hyperref}
\setlength{\parindent}{0pt}
\setlength{\parskip}{5pt}
\setlength{\emergencystretch}{3em}
\newcommand{\N}{\mathbb{N}}
\newcommand{\Z}{\mathbb{Z}}
\newcommand{\Q}{\mathbb{Q}}
\newcommand{\R}{\mathbb{R}}
\newcommand{\C}{\mathbb{C}}
\newcommand{\F}{\mathbb{F}}
\newcommand{\A}{\mathbb{A}}
\newcommand{\PP}{\mathbb{P}}
\newcommand{\E}{\mathbb{E}}
\newcommand{\eps}{\varepsilon}
\newcommand{\dd}{\,\mathrm{d}}
\begin{document}
\section*{UnsolvedMath ID: 10400010}
\textbf{Source problem code:} catalogue entry
\subsection{Definitions and mathematical statement}
Let \(K\subset S^3\) be a knot whose complement admits a complete finite-volume hyperbolic metric of sectional curvature \(-1\), and let \(\operatorname{Vol}(K)\) denote that volume. Let \(V_K(t)\) be the normalized Jones Laurent polynomial, with unknot value one. Set
\[
D_K=|V_K(-1)|,
\qquad A_K=\log D_K.
\]
The absolute value agrees with the supplied research summary's reading and avoids taking the real logarithm of a possibly negative Jones evaluation. In the usual knot normalization \(D_K\) is the determinant.

Dunfield asks to determine the relationship of \(\operatorname{Vol}(K)\) to \(A_K\), and separately to
\[
B_K=\frac{\log D_K}{\log d_K},\qquad d_K=\deg V_K(t)
\]
in the intended degree convention and where the denominator is defined and nonzero. For a Laurent polynomial, largest exponent and breadth
\[
\max\{j:[t^j]V_K\ne0\}-\min\{j:[t^j]V_K\ne0\}
\]
are different degree conventions; mirror image can also affect the first. A real logarithm requires \(d_K>0\) and \(d_K\ne1\). Breadth is a possible intended repair, but is not silently substituted for the source's degree.

The source supplies an empirical comparison of \(\operatorname{Vol}(K)\) and \(\pi\log D_K\) for alternating knots with 13 crossings, referring to Figure 4 and [112]. It does not specify a universal equation, inequalities, asymptotic regime, or knot class for the full two-part relationship question. The plotted ticks and dataset text are figure material, not extra numerical hypotheses.

The supplied research summary reports numerical partial understanding and mentions the Volume Conjecture for colored Jones asymptotics. That concerns a different family/evaluation limit and does not itself prove a relationship for the single Jones value at \(-1\), or for the degree-normalized expression here.
\subsection{Short English statement}
Relate hyperbolic knot volume to \(\log|V_K(-1)|\) and to its stated degree-normalized version. The source provides a 13-crossing alternating-knot plot, but no precise universal bound; sign and Laurent-degree conventions need clarification.
\subsection{Sources}
\begin{itemize}
\item Local UnsolvedMath registry, permanent ID \(10400010\), exact source code \texttt{catalogue entry}, original Problem 1.10, attributed to N. Dunfield.
\item Ohtsuki, \emph{Problems on invariants of knots and 3-manifolds} (2002), Problem 1.10, recorded PDF page 15; \url{https://msp.org/gtm/2002/04/gtm-2002-04-024s.pdf}. Recorded access 2026-07-29; pdf-font-structure extraction.
\item Figure 4, alternating knots with 13 crossings, dataset label \texttt{13\_alt.data}, source reference [112]; the bibliography entry is absent.
\item Supplied PARTIAL-PROGRESS research summary mentions the colored-Jones Volume Conjecture without a complete citation or theorem for the fixed-evaluation relationship.
\end{itemize}
\end{document}
